\section{Conclusion}
\label{sec:conclusion}

We set out to measure whether ML/NLP and HCI alignment researchers cite each other symmetrically within the huashen218 bidirectional alignment corpus.
We found that the measurement infrastructure was missing --- and built it.
We then discovered that the ``corpus'' available through the public GitHub reading list contains 49 extractable paper identifiers, not the $\sim$400 papers in the full systematic review.
In the 33 papers we could resolve, the asymmetry we hypothesized was already visible: a directed citation ratio of 0.107, with HCI alignment papers allocating 100\% of their within-corpus outgoing citations to ML/NLP work.
The full statistical test awaits a larger corpus. The pipeline to run it is validated and ready.

We contribute four things: a validated S2AG bibliometric pipeline (23/23 unit tests, fully cached); FoS-primary classification achieving 100\% vs.\ 12\% label coverage on interdisciplinary alignment corpora; the first directional citation measurement in the huashen218 corpus (ratio $\approx$ 0.107, directionally consistent with H-CitAsym but not statistically confirmable at $N=9$); and a documented methodological finding that curated reading lists yield only $\sim$12\% of the described corpus size --- programmatic S2AG citation-of-citation search is the viable reconstruction approach.

The most immediate extension is corpus augmentation (h-e1-v2): using the S2AG endpoint for papers citing \citet{Shen2024Position} to recover 200+ papers and enable full statistical testing.
All pipeline components are already implemented and validated; corpus augmentation is the remaining step.
From the unverified assumptions, temporal analysis (pre-/post-2022 citation behavior) becomes tractable with a larger corpus, as does venue string normalization enabling the pre-registered 3-scheme robustness analysis.

The measurement of directed citation asymmetry within the alignment community is the first step toward a structural account of whether ``bidirectional alignment'' is operationally bidirectional.
We hope this work provides both the methodology and the motivation for completing that measurement.
